% mage-auto-heading: the-capability
% mage-auto-heading: when-this-stack-earns-its-keep
% mage-auto-heading: the-composition
% mage-auto-heading: constituent-moves
% mage-auto-heading: why-these-travel-together
% Authoritative LaTeX manuscript; Markdown is a generated compatibility projection.

\section{The Capability}\label{heading:the-capability}

\textbf{Detect when the running system diverges from its operating model and trigger a bounded response.}

\section{When This Stack Earns Its Keep}\label{heading:when-this-stack-earns-its-keep}

Reach for it when:

\begin{itemize}
\tightlist
\item
  \textbf{The system runs live}, and a divergence between expected and observed state has real cost.
\item
  \textbf{Detection and response have drifted apart} --- a dashboard nobody acts on, or a runbook nobody triggers.
\item
  \textbf{Continuing can be more dangerous than stopping}, so an unsafe state needs a mechanism empowered to halt progress, not just a note in a log.
\end{itemize}

\section{The Composition}\label{heading:the-composition}

\begin{figure}[htbp]
\centering
\includegraphics[width=\linewidth]{book/latex/assets/observe-react-stack-6f07191f66.pdf}
\caption{The observe → react composition. An operating model drives WATCH, which reads runtime state; WATCH flows to RESPOND, which recovers. That solid path is required. Two dashed attachments strengthen it where the failure class demands: BEAT adds a liveness heartbeat so a hung process reads differently from a slow one, and BLOCK adds the power to prevent unsafe progress. Solid path: the core path; dashed: enhancement. Dashed attachment: a useful enhancement, not required for the capability.}
\label{observe-react-stack}
\end{figure}

\section{Constituent Moves}\label{heading:constituent-moves}

{\footnotesize
{\def\LTcaptype{none} % do not increment counter
\begin{longtable}[]{@{}
  >{\raggedright\arraybackslash}p{(\linewidth - 2\tabcolsep) * \real{0.5000}}
  >{\raggedright\arraybackslash}p{(\linewidth - 2\tabcolsep) * \real{0.5000}}@{}}
\toprule\noalign{}
\begin{minipage}[b]{\linewidth}\raggedright
Move
\end{minipage} & \begin{minipage}[b]{\linewidth}\raggedright
Role
\end{minipage} \\
\midrule\noalign{}
\endhead
\bottomrule\noalign{}
\endlastfoot
\textbf{WATCH} & Read runtime state against the operating model, so divergence surfaces as a signal. \\
\textbf{RESPOND} & Connect each signal to a bounded recovery action. \\
\textbf{BEAT} & \emph{(strengthens)} Emit periodic liveness so a wedged process reads differently from a slow one. \\
\textbf{BLOCK} & \emph{(strengthens)} Refuse unsafe progress where continuing is worse than stopping. \\
\end{longtable}
}
}

\section{Why These Travel Together}\label{heading:why-these-travel-together}

Observation without response is a dashboard. Response without observation is a runbook waiting for someone to notice the fault.

WATCH and RESPOND form the core loop: observe divergence from the expected state, then trigger a bounded response. Add a heartbeat when silence must be distinguished from slow progress; add a blocking mechanism when continuing is more dangerous than stopping.

\textbf{Mechanisms:} watchdog · heartbeat · bounded recovery · fail-closed gate
